\section{Scope, results, and the source questions}
\label{sec:scope}

The source corpus is already extensive. It contains general harmonic
interpolation, centered tail moments, Gamma and Stieltjes derivative
ladders, finite spectral-subtraction algorithms, and a careful record
of rejected or still-conjectural period relations. A continuation should
therefore do more than rederive a familiar Euler sum or enlarge a
numerical search. We select explicit unresolved questions, identify
their inherited inputs, and supply complete new proof steps.

The baseline is commit
\texttt{da2792fc9cc8ef19c6a314be6c65e1437fbdc66b}, including the canonical
manuscript \cite{RepoMain}, the harmonic-gap continuation
\cite{gm:gaps-source}, and all five incoming archives dated 11 October
2026. In particular the quartic reconstruction report \cite{RepoQuartic}
already asks for fifth- and higher-power cancellation, the complementary
resonant parity, and completion of the order-ten reconstruction family.
The higher-defect report \cite{sd:HigherDefects} asks for explicit
simplex-face subtractions and their odd-weight extension. These are
concrete targets addressed here.

\subsection{The main new statements}

The first contribution is an all-power mechanism. If
\[
 A_n=\psi'(n+1)=\zeta(2)-H_n^{(2)},\qquad x_n=n+\tfrac12,
\]
the regular centered moments of $A_n^d$ have an anchored generating
function satisfying one triangular recurrence for every $d$. Its finite
boundary cancellation is the parity of
\[
 u^{2d}B_d\!\left(\tfrac12+\frac{E(u)}u\right),
 \qquad E(u)=\sum_{j\ge0}B_{2j}u^{2j}.
\]
This polynomial identity explains the lower-power examples and proves
the fixed finite-span conclusion without extrapolating from a finite
table. The resulting fifth-power coordinate system has eight formal
coordinates, and the sixth-power system has eleven. Exact rational
determinants certify the full ranks of the displayed coefficient rows.
An especially compact consequence is
\begin{equation}\label{intro:quintic}
\begin{split}
&\sum_{n=0}^{\infty}\Bigl[
 (195x_n+4056x_n^3+11568x_n^5+5376x_n^7)A_n^5
                  -5376x_n^2-9328\Bigr]\\
&\hspace{12mm}=-748\zeta(2)+7884\zeta(3)-5380\zeta(4)
                                    +171\zeta(6)-748.
\end{split}
\end{equation}
The entire bracket is one summand. Its series is ordinary and absolutely
convergent; splitting off its divergent polynomial pieces would change
the meaning. The proof is finite summation and elimination, not integer
relation recognition.

The same mechanism works for every repeated even-order Hurwitz tail
\(\zeta(2q,n+1)^d\), with an explicit finite span of weights at most
$2qd-1$ (Theorem~\ref{et:main}). A short consequence, with
$T_n=\psi^{(3)}(n+1)/6$, is
\[
 \sum_{n\ge0}\bigl((4x_n^3+x_n)T_n^3-3T_n^2\bigr)
                                      =-\frac94\zeta(8).
\]
The underlying cubic cutoff identity also gives a complex-order family
for $\Re p>3/2$. At the critical real order $p=3/2$, the ordinary
sum acquires the exact additional endpoint contribution $8$;
see Corollary~\ref{et:complex-cubic}.

The same spectral family has actual poles at $1-d-2r$. We give all their
Laurent coefficients as ordinary logarithmic sums plus a finite Hurwitz
and generalized-Stieltjes correction. Removing the pole gives a Taylor
series of exact radius two. This answers the normalization part of the
resonant question and supplies an analytic local generator, while leaving
open a uniform arithmetic reduction as the resonance index increases.

The second contribution completes a previously partial theorem. For
each even larger order $2b>10$, the first $2b$ regular centered moments
of $\zeta(10,n+1)\zeta(2b,n+1)$ have the optimal formal coordinate
rank $2b$. Two existing reports cover complementary ranges: a congruence
covers all but one residue class, and an eventual theorem covers
$b\ge1920$. Their only uncovered index is $b=962$. A five-by-five
matrix modulo $2^{20}$ proves nonvanishing there. We replay the inherited
finite polynomial identities and give an explicit rigorous large-order
comparison, so the completed family has no unchecked numerical case.
The additional parameter pair relative to the union of the two reports
is exactly $(10,1924)$.

The third contribution treats several transverse marks in relative
multiple-Hurwitz interpolation. Its first useful case is
\begin{equation}\label{intro:gamma}
\left.\partial_u\partial_v\mathcal H_{a,b}(s,u,v,t)\right|_{u=v=0}
=\frac12\mathcal E_{a,b}^{s,t}\!\left[
 (\log\Gamma(z+1)-\log\Gamma(y))^2
                  -2\int_y^{z+1}\gamma_1(x)\,dx\right].
\end{equation}
Here $\mathcal H$ is the prescribed relative quotient and $\mathcal E$
is an independently defined decorated double quotient; their full
definitions and ordinary-series realizations appear in
Section~\ref{gm:section}. The identity holds for arbitrary independent
spectral arguments and endpoints in the right half-plane. Its extension
uses a finite set-partition polynomial. Identical marked clusters give
individual ordered identities; different marks give their complete
permutation-symmetric sum. The analytic proof establishes decoration
closure before using the partition identity.

The fourth contribution resolves the face-subtraction problem for
higher spectral defects. For a polynomial weight, any number of
nonnegative masses, and every derivative order, we give an exact finite
sum of ordinary integrals over simplex faces. A $k$-vertex face first
appears at derivative order $k$, after the zero-jet part is separated.
That zero-jet part is an explicit rational function. It is absent for
even polynomial weights and, more generally, for every polynomial
weight multiplied by a mean-zero periodic sequence. The highest faces
at each derivative order require only the first one-mass derivative,
which for constant weight is an ordinary log-Gamma expression.
The subtraction also gives a closed identity at every face size:
\begin{equation}\label{intro:face}
 \int_{\Sigma_S}\mathcal H_S(B\mapsto B^2;t)
                  \frac{dt}{\prod_{i\in S}t_i}
 =(-1)^{|S|-1}(|S|-2)!
                 \sum_{\{i,j\}\subseteq S}(B_i-B_j)^2.
\end{equation}
The simplex normalization and face operator are specified in
Section~\ref{sd:sec:main}; the integral has no hidden finite part.

\subsection{Claim boundaries and relation to earlier work}

Several ingredients are classical and are used with attribution.
Hoffman's ordering of the minimum index gives products of zeta tails
as multiple zeta values \cite{HoffmanTails}. His symmetric-sum theorem
supplies the partition coefficients \cite{gm:Hoffman1992,gm:Hoffman2018}.
Ihara--Nakamura--Yamamoto supply the bare entire interpolation and its
harmonic product \cite{gm:INY}. The scalar Dirichlet mixture is a
classical Gamma-integral identity, already present in the source; the
literature on logarithmic Dirichlet averages includes
\cite{sd:Carlson1987,sd:Carlson1991}. These facts are inputs, not claimed
discoveries of the present article.

The new content is the all-power boundary mechanism and its concrete
consequences; the missing exact reconstruction certificate and completed
family; the local analytic realization of multiple decorations and
marks; and the explicit face decomposition with its rational correction
and all-order filtration. Novelty is assessed relative to the inspected
repository and the cited primary literature. An exhaustive priority
claim over the entire literature is not made.

The manuscript's Gaussian $S_6$ and revised $S_8$ period identities
remain open. The rejected historical $S_8$ vector is a different object
and retains its rejected status. No determinant of a formal coordinate
matrix proves arithmetic independence of evaluated periods. Likewise a
representation by polylogarithms of bounded depth is an upper bound in
that representation, not a proof of minimal depth. These distinctions
are part of the theorem statements, not qualifications added only after
the calculations.

\subsection{Conventions and proof architecture}

We use $H_n^{(r)}=\sum_{j=1}^n j^{-r}$, with $H_0^{(r)}=0$;
Bernoulli numbers have $B_1=-1/2$; and multiple zeta values are ordered
as
\[
 \zeta(s_1,\ldots,s_\ell)
 =\sum_{n_1>\cdots>n_\ell\ge1}
                         n_1^{-s_1}\cdots n_\ell^{-s_\ell}.
\]
The prime on a zeta function denotes differentiation in its spectral
argument unless explicitly stated otherwise. The generalized Stieltjes
constants satisfy
\[
 \zeta(1+\varepsilon,a)=\frac1\varepsilon+
       \sum_{q\ge0}\frac{(-1)^q\gamma_q(a)}{q!}\varepsilon^q.
\]
Section-specific letters are reset explicitly: $F_d(s)$ denotes a tail
spectrum, whereas $F_j(B)$ in the spectral-face section denotes a
one-mass derivative. Every antiderivative has its endpoint fixed.

The arguments use three distinct kinds of evidence. Finite identities
are proved algebraically and supplied with rational or modular replay.
Analytic statements include convergence domains, continuation from an
open domain, and endpoint normalization. High-precision evaluations
then check independent representations and signs. Numerical agreement
is not used as a substitute for either of the first two steps.
Section~\ref{sec:audit} records what was audited and what the artifact
checks establish; Section~\ref{sec:research} lists exact next targets.
